\uuid{O7yO}
\exo7id{834}
\titre{exo7 834}
\auteur{ridde}
\organisation{exo7}
\datecreate{1999-11-01}
\isIndication{false}
\isCorrection{false}
\chapitre{Calcul d'intégrales}
\sousChapitre{Primitives diverses}
\cpge{mpsi-mp2i.primitives-equations-differentielles.a, pcsi.techniques-calcul-primitives.a, ptsi.techniques-calcul-primitives.a, mpsi-mp2i.fonctions-valeurs-reelles-complexes.a, pcsi.techniques-calcul-fonctions.a, ptsi.techniques-calcul-fonctions.a, mpsi-mp2i.polynomes-fractions-rationnelles.h, bcpst1.derivees-primitives-integrales.c, bcpst1.derivees-primitives-integrales.d, bcpst1.fonctions-reelles-usuelles.a}
\module{Analyse}
\niveau{L1}
\difficulte{}

\contenu{
\question{
D\'eterminer les intervalles d'\'etude et calculer les primitives des fonctions :
$$ \text{ch}x \sin (2x)$$
$$ \frac 1{\sqrt{2 + \tan ^2 x}}$$
$$ (x^2 + 2x + 2) \cos (2x)$$
$$ x^2 \cos x \text{ et } x^2 \sin x \text{ en utilisant les complexes}$$
$$\frac 1 { (x^2-1)^3} \text{ et } \frac 1{ (x^2-1)^2}$$
$$ \frac {\sqrt{1 + x}} {x \sqrt{1-x}}$$
}
}
